\begin{ExerciseBlock}
\ExerciseHeading{05}{Stehende Welle mit offenem und geschlossenem Ende}{doppler_and_interference_pan_flute}
Das Rohr einer Panflöte ist an einem Ende geschlossen und am anderen offen. An beiden Enden werden die Schallwellen reflektiert. Die Rohrlänge ist $L=\qty{20}{\centi\metre}$, die Schallgeschwindigkeit beträgt $c_\mathrm L\approx\qty{340}{\metre\per\second}$. Vernachlässigen Sie die Mündungskorrektur.
\begin{center}
\CachedTikzFigure{content/resources/chapters/doppler_and_interference/tikz/pan_flute.tex}
\end{center}
\begin{enumerate}
\item Welche Randbedingungen gelten für die ortsabhängige Auslenkungsamplitude $s_0(x)$ an beiden Rohrenden?
\item Skizzieren Sie $s_0(x)$ für die Grundmode und die erste Oberschwingung.
\item Verwenden Sie den Ansatz $s_0(x)=A\sin(kx)+B\cos(kx)$. Bestimmen Sie die Bedingungen an $A$, $B$ und $k$, die beide Randbedingungen erfüllen. Geben Sie alle möglichen Wellenzahlen $k_n$ und ihre Werte für die Grundmode und die erste Oberschwingung an.
\item Bestimmen Sie allgemein die Frequenzen $f_n$ und Wellenlängen $\lambda_n$ sowie deren Werte für die beiden genannten Moden.
\end{enumerate}
\end{ExerciseBlock}
\UseExerciseSolution{doppler_and_interference_pan_flute}
